\subsection{Bisect $k$-Adjustment for MMDs}

The bisect algorithm generates a $k$-partition of a state's census tracts
into equal-population districts.  For single-member redistricting, $k$ equals
the number of congressional seats apportioned to the state.  For MMDs, $k$
is reduced to the number of districts, which is $S / M$ for a uniform-magnitude
design (where $S$ is total seats and $M$ is seats per district):

\begin{equation}
  k_{\text{MMD}} = \left\lfloor \frac{S}{M} \right\rfloor
\end{equation}

For non-divisible cases (e.g., Illinois with $S = 17$, $M = 3$: $k = 5$
with one district taking 4 seats and four districts taking 3), we implement
a mixed-magnitude design where the algorithm generates 5 equal-population
partitions and the seat count per district is assigned post-hoc based on
population ranking.

The bisect CLI exposes this through the \texttt{--seats-per-district} flag:

\begin{verbatim}
bisect state --state CA --year 2020 --version E1_mmd4 \
  --districts 13 --seats-per-district 4 \
  --balance-tolerance 5.0 --seed 42
\end{verbatim}

\subsection{Per-District Balance Tolerance}

A critical algorithmic adjustment for MMDs is the balance tolerance.
Single-member redistricting targets $\pm 0.5\%$ population deviation from
the ideal district population ($P / k$ where $P$ is state population).  This
strict tolerance reflects the ``one person, one vote'' constitutional
standard for congressional districts \citep{wesberry1964}.

For MMDs, the relevant equality standard is \emph{per seat}, not per
district: equal representation means each seat represents the same number of
people, regardless of how seats are grouped into districts.  For a $k$-district
design with $M$ seats per district:

\begin{equation}
  \text{ideal per seat} = \frac{P}{k \cdot M} = \frac{P}{S}
\end{equation}

\begin{equation}
  \text{ideal per district} = M \cdot \frac{P}{S}
\end{equation}

The tolerance per district is therefore $M$ times the per-seat tolerance.
If we target $\pm 0.5\%$ per seat, the per-district tolerance is
$\pm 0.5\% \cdot M$.  For 4-member districts this is $\pm 2.0\%$;
for 2-member districts, $\pm 1.0\%$.

However, this relaxed tolerance is appropriate for the population balance
constraint only if the \emph{per-seat} standard is applied when evaluating
constitutional compliance.  We implement both: the algorithm uses the
per-district tolerance for feasibility, and we report per-seat deviations
in our results.

\subsection{METIS Ufactor Scaling}

The METIS partitioner uses an ``imbalance factor'' (ufactor) to control
population balance at each bisection node.  As described in prior portfolio
work \citep{dellib2026edgeweighted}, the ufactor must be scaled by district
count to prevent compounding balance errors at each bisection level.

For MMD generation, the ufactor formula is:
\begin{equation}
  \text{ufactor}_{\text{node}} = 1 + \frac{T \cdot M}{k_{\text{node}}}
\end{equation}
where $T$ is the per-seat tolerance (0.005 for $\pm 0.5\%$), $M$ is seats
per district, and $k_{\text{node}}$ is the number of districts in the current
subtree.  This guarantees cumulative balance error $\leq T \cdot M$ regardless
of bisection depth.

\subsection{STV Implementation}

We implement the Gregory method of STV \citep{gallagher2011ireland}, which
uses fractional surplus transfer:

\begin{enumerate}
  \item Compute Droop quota $Q = \lfloor P / (S+1) \rfloor + 1$.
  \item Award seats to all candidates with votes $\geq Q$.  Transfer surplus
        at value $TV = (v_c - Q) / v_c$ applied to all continuing ballots
        for that candidate.
  \item If no candidate reaches $Q$, eliminate the candidate with fewest
        first-preference votes.  Transfer their votes at full value.
  \item Repeat until all $S$ seats are filled or remaining candidates equal
        remaining vacancies.
\end{enumerate}

For simulation purposes (where we do not have individual ranked ballots), we
use vote-share input: parties receive their seat-proportion of the Droop quota
from aggregate vote data.  This approximation is standard in comparative
electoral studies \citep{gallagher1991proportionality} and produces accurate
seat counts when parties are cohesive.

\subsection{OLPR Implementation}

OLPR seat allocation uses the D'Hondt highest-averages method:

\begin{enumerate}
  \item For each party $p$ with vote total $v_p$, compute quotients
        $v_p / (s_p + 1)$ where $s_p$ is the current seat count.
  \item Award the next seat to the party with the highest quotient.
  \item Repeat until all $S$ seats are awarded.
\end{enumerate}

The effective threshold---the minimum vote share needed to win at least one
seat---is approximately $1/(S+1)$ (Droop quota): 20\% for 4-member districts,
25\% for 3-member, 33\% for 2-member.

\subsection{Data and Configurations}

We run three state case studies using 2020 census tract data and 2020
presidential election returns \citep{MEDSL2020}:

\begin{description}
  \item[California ($S=52$).] Two configurations: $k=13$ districts of 4
    seats (``CA-4''), and $k=26$ districts of 2 seats (``CA-2'').
  \item[New Jersey ($S=12$).] Two configurations: $k=4$ districts of 3 seats
    (``NJ-3''), and $k=6$ districts of 2 seats (``NJ-2'').
  \item[Illinois ($S=17$).] One configuration: $k=5$ districts with a
    mixed magnitude (three 3-seat and two 4-seat districts, ``IL-mix'').
\end{description}

All runs use seed 42, geographic edge-weighting, and the \texttt{single}
search mode.  Results carry the single-run dagger (\singrun) throughout.
